% Options for packages loaded elsewhere
\PassOptionsToPackage{unicode}{hyperref}
\PassOptionsToPackage{hyphens}{url}
\documentclass[
  a4paper,
]{ltjsarticle}
\usepackage{xcolor}
\usepackage[margin=25mm]{geometry}
\usepackage{amsmath,amssymb}
\setcounter{secnumdepth}{-\maxdimen} % remove section numbering
\usepackage{iftex}
\ifPDFTeX
  \usepackage[T1]{fontenc}
  \usepackage[utf8]{inputenc}
  \usepackage{textcomp} % provide euro and other symbols
\else % if luatex or xetex
  \usepackage{unicode-math} % this also loads fontspec
  \defaultfontfeatures{Scale=MatchLowercase}
  \defaultfontfeatures[\rmfamily]{Ligatures=TeX,Scale=1}
\fi
\usepackage{lmodern}
\ifPDFTeX\else
  % xetex/luatex font selection
\fi
% Use upquote if available, for straight quotes in verbatim environments
\IfFileExists{upquote.sty}{\usepackage{upquote}}{}
\IfFileExists{microtype.sty}{% use microtype if available
  \usepackage[]{microtype}
  \UseMicrotypeSet[protrusion]{basicmath} % disable protrusion for tt fonts
}{}
\makeatletter
\@ifundefined{KOMAClassName}{% if non-KOMA class
  \IfFileExists{parskip.sty}{%
    \usepackage{parskip}
  }{% else
    \setlength{\parindent}{0pt}
    \setlength{\parskip}{6pt plus 2pt minus 1pt}}
}{% if KOMA class
  \KOMAoptions{parskip=half}}
\makeatother
\usepackage{color}
\usepackage{fancyvrb}
\newcommand{\VerbBar}{|}
\newcommand{\VERB}{\Verb[commandchars=\\\{\}]}
\DefineVerbatimEnvironment{Highlighting}{Verbatim}{commandchars=\\\{\}}
% Add ',fontsize=\small' for more characters per line
\newenvironment{Shaded}{}{}
\newcommand{\AlertTok}[1]{\textcolor[rgb]{1.00,0.00,0.00}{\textbf{#1}}}
\newcommand{\AnnotationTok}[1]{\textcolor[rgb]{0.38,0.63,0.69}{\textbf{\textit{#1}}}}
\newcommand{\AttributeTok}[1]{\textcolor[rgb]{0.49,0.56,0.16}{#1}}
\newcommand{\BaseNTok}[1]{\textcolor[rgb]{0.25,0.63,0.44}{#1}}
\newcommand{\BuiltInTok}[1]{\textcolor[rgb]{0.00,0.50,0.00}{#1}}
\newcommand{\CharTok}[1]{\textcolor[rgb]{0.25,0.44,0.63}{#1}}
\newcommand{\CommentTok}[1]{\textcolor[rgb]{0.38,0.63,0.69}{\textit{#1}}}
\newcommand{\CommentVarTok}[1]{\textcolor[rgb]{0.38,0.63,0.69}{\textbf{\textit{#1}}}}
\newcommand{\ConstantTok}[1]{\textcolor[rgb]{0.53,0.00,0.00}{#1}}
\newcommand{\ControlFlowTok}[1]{\textcolor[rgb]{0.00,0.44,0.13}{\textbf{#1}}}
\newcommand{\DataTypeTok}[1]{\textcolor[rgb]{0.56,0.13,0.00}{#1}}
\newcommand{\DecValTok}[1]{\textcolor[rgb]{0.25,0.63,0.44}{#1}}
\newcommand{\DocumentationTok}[1]{\textcolor[rgb]{0.73,0.13,0.13}{\textit{#1}}}
\newcommand{\ErrorTok}[1]{\textcolor[rgb]{1.00,0.00,0.00}{\textbf{#1}}}
\newcommand{\ExtensionTok}[1]{#1}
\newcommand{\FloatTok}[1]{\textcolor[rgb]{0.25,0.63,0.44}{#1}}
\newcommand{\FunctionTok}[1]{\textcolor[rgb]{0.02,0.16,0.49}{#1}}
\newcommand{\ImportTok}[1]{\textcolor[rgb]{0.00,0.50,0.00}{\textbf{#1}}}
\newcommand{\InformationTok}[1]{\textcolor[rgb]{0.38,0.63,0.69}{\textbf{\textit{#1}}}}
\newcommand{\KeywordTok}[1]{\textcolor[rgb]{0.00,0.44,0.13}{\textbf{#1}}}
\newcommand{\NormalTok}[1]{#1}
\newcommand{\OperatorTok}[1]{\textcolor[rgb]{0.40,0.40,0.40}{#1}}
\newcommand{\OtherTok}[1]{\textcolor[rgb]{0.00,0.44,0.13}{#1}}
\newcommand{\PreprocessorTok}[1]{\textcolor[rgb]{0.74,0.48,0.00}{#1}}
\newcommand{\RegionMarkerTok}[1]{#1}
\newcommand{\SpecialCharTok}[1]{\textcolor[rgb]{0.25,0.44,0.63}{#1}}
\newcommand{\SpecialStringTok}[1]{\textcolor[rgb]{0.73,0.40,0.53}{#1}}
\newcommand{\StringTok}[1]{\textcolor[rgb]{0.25,0.44,0.63}{#1}}
\newcommand{\VariableTok}[1]{\textcolor[rgb]{0.10,0.09,0.49}{#1}}
\newcommand{\VerbatimStringTok}[1]{\textcolor[rgb]{0.25,0.44,0.63}{#1}}
\newcommand{\WarningTok}[1]{\textcolor[rgb]{0.38,0.63,0.69}{\textbf{\textit{#1}}}}
\setlength{\emergencystretch}{3em} % prevent overfull lines
\providecommand{\tightlist}{%
  \setlength{\itemsep}{0pt}\setlength{\parskip}{0pt}}
\usepackage{bookmark}
\IfFileExists{xurl.sty}{\usepackage{xurl}}{} % add URL line breaks if available
\urlstyle{same}
\hypersetup{
  hidelinks,
  pdfcreator={LaTeX via pandoc}}

\author{}
\date{}

\begin{document}

\section{無名等振幅複合波モデル基本公理系
v4}\label{ux7121ux540dux7b49ux632fux5e45ux8907ux5408ux6ce2ux30e2ux30c7ux30ebux57faux672cux516cux7406ux7cfb-v4}

\textbf{日付:} 2026-07-12\\
\textbf{著者:} 木原 範昭\\
\textbf{位置づけ:} 定義論文\\
\textbf{Version DOI:} 10.5281/zenodo.21316620 \textbf{Concept DOI:}
10.5281/zenodo.21315735

\begin{center}\rule{0.5\linewidth}{0.5pt}\end{center}

\section{1. 第一原理}\label{ux7b2cux4e00ux539fux7406}

\subsection{公理0:
成分無名性}\label{ux516cux74060-ux6210ux5206ux7121ux540dux6027}

分類: 公理

基本成分に、個別名、特権軸、特権符号、特権型を与えない。

添字 \texttt{n} は便宜上のラベルであり、成分の固有名ではない。

禁止:

\begin{Shaded}
\begin{Highlighting}[]
\NormalTok{特定成分だけを負符号軸にする}
\NormalTok{特定成分だけを時間軸にする}
\NormalTok{特定成分だけを実数型にする}
\NormalTok{特定成分だけを複素型にする}
\NormalTok{特定成分だけに固有名を与える}
\end{Highlighting}
\end{Shaded}

\begin{center}\rule{0.5\linewidth}{0.5pt}\end{center}

\subsection{公理0+:
定式化無名性}\label{ux516cux74060-ux5b9aux5f0fux5316ux7121ux540dux6027}

分類: 公理

第一原理層では、理論名、定式名、読出し論理名にも固有名を与えない。

物理名、対象名、相互作用名、既存理論上の分類名を理由として、第一原理層の方程式、射影、読出し規則を変更してはならない。

禁止:

\begin{Shaded}
\begin{Highlighting}[]
\NormalTok{重力だから円運動型の式を採用する}
\NormalTok{クーロン力だから双曲線型の式を採用する}
\NormalTok{質量だから符号を消す}
\NormalTok{電荷だから符号を残す}
\NormalTok{エネルギーだから時間軸へ移す}
\NormalTok{二体系だから別の閉鎖式を使う}
\NormalTok{三体系だから別の閉鎖式を使う}
\NormalTok{ABC系だから特別な読出し規則を使う}
\NormalTok{欲しい物理名に合わせて右辺表示を選ぶ}
\end{Highlighting}
\end{Shaded}

許可:

\begin{Shaded}
\begin{Highlighting}[]
\NormalTok{同一閉鎖条件に対する読出し演算の違い}
\NormalTok{同一閉鎖条件に対する対称性の破れの違い}
\NormalTok{同一閉鎖条件に対するゲージ安定性の違い}
\NormalTok{同一閉鎖条件に対する符号保持の違い}
\NormalTok{同一閉鎖条件に対する符号消失の違い}
\NormalTok{同一閉鎖条件に対する位相枝選択の違い}
\NormalTok{同一閉鎖条件に対する閉鎖対象集合の違い}
\end{Highlighting}
\end{Shaded}

順序:

\begin{Shaded}
\begin{Highlighting}[]
\NormalTok{同一の全正符号ゼロ閉鎖を置く}
\NormalTok{読出し演算を先に定義する}
\NormalTok{対称性の破れを先に定義する}
\NormalTok{符号保持または符号消失を先に定義する}
\NormalTok{ゲージ安定性を確認する}
\NormalTok{読出し後に作業名を与える}
\end{Highlighting}
\end{Shaded}

対象集合:

\begin{Shaded}
\begin{Highlighting}[]
\NormalTok{Q(P)=\textbackslash{}sum\_n p\_n\^{}2}
\end{Highlighting}
\end{Shaded}

\begin{Shaded}
\begin{Highlighting}[]
\NormalTok{Q(A)=0}
\end{Highlighting}
\end{Shaded}

\begin{Shaded}
\begin{Highlighting}[]
\NormalTok{Q(A+B)=0}
\end{Highlighting}
\end{Shaded}

\begin{Shaded}
\begin{Highlighting}[]
\NormalTok{Q(A+B+C)=0}
\end{Highlighting}
\end{Shaded}

対象集合の違いは許される。

\texttt{Q} の定義変更は許されない。

\begin{center}\rule{0.5\linewidth}{0.5pt}\end{center}

\subsection{公理0.5:
次元別中心投影読出し}\label{ux516cux74060.5-ux6b21ux5143ux5225ux4e2dux5fc3ux6295ux5f71ux8aadux51faux3057}

分類: 補助幾何公理

単一軸上の読出しでは曲率補正を導入しない。

曲率補正候補は、独立な二方向以上が面積、体積、閉路、または輸送不一致を形成する場合に限り導入する。

\begin{center}\rule{0.5\linewidth}{0.5pt}\end{center}

\subsection{公理1:
全正符号ゼロ閉鎖}\label{ux516cux74061-ux5168ux6b63ux7b26ux53f7ux30bcux30edux9589ux9396}

分類: 公理

\begin{Shaded}
\begin{Highlighting}[]
\NormalTok{\textbackslash{}sum\_\{n=1\}\^{}\{N\}x\_n\^{}2=0}
\end{Highlighting}
\end{Shaded}

展開形:

\begin{Shaded}
\begin{Highlighting}[]
\NormalTok{x\_1\^{}2+x\_2\^{}2+\textbackslash{}cdots+x\_N\^{}2=0}
\end{Highlighting}
\end{Shaded}

係数符号:

\begin{Shaded}
\begin{Highlighting}[]
\NormalTok{全項 +}
\end{Highlighting}
\end{Shaded}

非採用:

\begin{Shaded}
\begin{Highlighting}[]
\NormalTok{\textbackslash{}sum\_n |x\_n|\^{}2=0}
\end{Highlighting}
\end{Shaded}

閉鎖に用いる項:

\begin{Shaded}
\begin{Highlighting}[]
\NormalTok{x\_n\^{}2}
\end{Highlighting}
\end{Shaded}

閉鎖に用いない項:

\begin{Shaded}
\begin{Highlighting}[]
\NormalTok{x\_n\textbackslash{}bar\{x\}\_n}
\end{Highlighting}
\end{Shaded}

適用状態:

\begin{Shaded}
\begin{Highlighting}[]
\NormalTok{定常閉鎖状態}
\NormalTok{安定閉鎖状態}
\end{Highlighting}
\end{Shaded}

準定常状態:

\begin{Shaded}
\begin{Highlighting}[]
\NormalTok{外部影響直後}
\NormalTok{閉鎖回復過程}
\NormalTok{準安定状態}
\end{Highlighting}
\end{Shaded}

準定常状態における一時的な非ゼロ閉鎖残差:

\begin{Shaded}
\begin{Highlighting}[]
\NormalTok{\textbackslash{}sum\_n x\_n\^{}2\textbackslash{}ne0}
\end{Highlighting}
\end{Shaded}

定常化後:

\begin{Shaded}
\begin{Highlighting}[]
\NormalTok{\textbackslash{}sum\_n x\_n\^{}2=0}
\end{Highlighting}
\end{Shaded}

\begin{center}\rule{0.5\linewidth}{0.5pt}\end{center}

\subsection{公理2:
非自明存在}\label{ux516cux74062-ux975eux81eaux660eux5b58ux5728}

分類: 公理

\begin{Shaded}
\begin{Highlighting}[]
\NormalTok{\textbackslash{}exists n,\textbackslash{}quad x\_n\textbackslash{}neq0}
\end{Highlighting}
\end{Shaded}

\begin{center}\rule{0.5\linewidth}{0.5pt}\end{center}

\subsection{公理3:
投影軸縮退と虚数方向読出し}\label{ux516cux74063-ux6295ux5f71ux8ef8ux7e2eux9000ux3068ux865aux6570ux65b9ux5411ux8aadux51faux3057}

分類: 公理

投影法により内在座標系が構成されるとき、投影軸方向は、その内在座標系の住民から方向として判別されない。

ただし、投影軸方向の値が閉鎖条件内に残る場合、その値を消去しない。

投影軸方向の補償量を \texttt{z} と読むとき、内在座標系ではこれを

\begin{Shaded}
\begin{Highlighting}[]
\NormalTok{iz}
\end{Highlighting}
\end{Shaded}

として扱う。

この成分は一次方向としては観測不能である。

閉鎖式内で二乗されると、

\begin{Shaded}
\begin{Highlighting}[]
\NormalTok{(iz)\^{}2={-}z\^{}2}
\end{Highlighting}
\end{Shaded}

となる。

したがって、投影軸方向の観測不能な補償量は、二乗読出しを通じて、符号反転した補償量として内在座標系に現れる。

\begin{center}\rule{0.5\linewidth}{0.5pt}\end{center}

\section{2.
第一原理からの帰結}\label{ux7b2cux4e00ux539fux7406ux304bux3089ux306eux5e30ux7d50}

\subsection{帰結1:
実数正定値成分の不足}\label{ux5e30ux7d501-ux5b9fux6570ux6b63ux5b9aux5024ux6210ux5206ux306eux4e0dux8db3}

分類: 導出済み帰結

\begin{Shaded}
\begin{Highlighting}[]
\NormalTok{x\_n\textbackslash{}in\textbackslash{}mathbb\{R\}}
\NormalTok{\textbackslash{}Rightarrow}
\NormalTok{x\_n\^{}2\textbackslash{}ge0}
\end{Highlighting}
\end{Shaded}

\begin{Shaded}
\begin{Highlighting}[]
\NormalTok{\textbackslash{}sum\_n x\_n\^{}2=0}
\NormalTok{\textbackslash{}Rightarrow}
\NormalTok{x\_n=0\textbackslash{}quad\textbackslash{}forall n}
\end{Highlighting}
\end{Shaded}

\begin{center}\rule{0.5\linewidth}{0.5pt}\end{center}

\subsection{帰結2:
符号反転の内部生成}\label{ux5e30ux7d502-ux7b26ux53f7ux53cdux8ee2ux306eux5185ux90e8ux751fux6210}

分類: 導出済み帰結

\begin{Shaded}
\begin{Highlighting}[]
\NormalTok{x\_1=A,\textbackslash{}qquad x\_2=iA}
\end{Highlighting}
\end{Shaded}

\begin{Shaded}
\begin{Highlighting}[]
\NormalTok{x\_1\^{}2+x\_2\^{}2=A\^{}2+(iA)\^{}2=A\^{}2{-}A\^{}2=0}
\end{Highlighting}
\end{Shaded}

\begin{Shaded}
\begin{Highlighting}[]
\NormalTok{i\^{}2={-}1}
\end{Highlighting}
\end{Shaded}

\begin{center}\rule{0.5\linewidth}{0.5pt}\end{center}

\subsection{帰結3:
位相代数の必要性}\label{ux5e30ux7d503-ux4f4dux76f8ux4ee3ux6570ux306eux5fc5ux8981ux6027}

分類: 導出済み帰結

\begin{Shaded}
\begin{Highlighting}[]
\NormalTok{x\_n=r\_n e\^{}\{i\textbackslash{}phi\_n\}}
\end{Highlighting}
\end{Shaded}

\begin{Shaded}
\begin{Highlighting}[]
\NormalTok{x\_n\^{}2=r\_n\^{}2e\^{}\{i2\textbackslash{}phi\_n\}}
\end{Highlighting}
\end{Shaded}

閉鎖条件:

\begin{Shaded}
\begin{Highlighting}[]
\NormalTok{\textbackslash{}sum\_\{n=1\}\^{}\{N\}r\_n\^{}2e\^{}\{i2\textbackslash{}phi\_n\}=0}
\end{Highlighting}
\end{Shaded}

等振幅条件:

\begin{Shaded}
\begin{Highlighting}[]
\NormalTok{r\_n=r}
\end{Highlighting}
\end{Shaded}

\begin{Shaded}
\begin{Highlighting}[]
\NormalTok{\textbackslash{}sum\_\{n=1\}\^{}\{N\}e\^{}\{i2\textbackslash{}phi\_n\}=0}
\end{Highlighting}
\end{Shaded}

\begin{center}\rule{0.5\linewidth}{0.5pt}\end{center}

\subsection{帰結4:
型無名性}\label{ux5e30ux7d504-ux578bux7121ux540dux6027}

分類: 導出済み帰結

\begin{Shaded}
\begin{Highlighting}[]
\NormalTok{x\_n\textbackslash{}in\textbackslash{}mathbb\{C\}\textbackslash{}quad\textbackslash{}forall n}
\end{Highlighting}
\end{Shaded}

または、

\begin{Shaded}
\begin{Highlighting}[]
\NormalTok{全成分に同一の位相代数型を与える}
\end{Highlighting}
\end{Shaded}

\begin{center}\rule{0.5\linewidth}{0.5pt}\end{center}

\subsection{帰結5:
90度位相差}\label{ux5e30ux7d505-90ux5ea6ux4f4dux76f8ux5dee}

分類: 導出済み帰結

\begin{Shaded}
\begin{Highlighting}[]
\NormalTok{A\^{}2+(iA)\^{}2=0}
\end{Highlighting}
\end{Shaded}

\section{3.
無名等振幅複合波}\label{ux7121ux540dux7b49ux632fux5e45ux8907ux5408ux6ce2}

\subsection{公理4:
無名等振幅複合波}\label{ux516cux74064-ux7121ux540dux7b49ux632fux5e45ux8907ux5408ux6ce2}

分類: 公理

\begin{Shaded}
\begin{Highlighting}[]
\NormalTok{\textbackslash{}Psi\_N=\textbackslash{}frac\{A\_0\}\{N\}\textbackslash{}sum\_\{k=1\}\^{}\{N\}e\^{}\{i\textbackslash{}theta\_k\}}
\end{Highlighting}
\end{Shaded}

\begin{Shaded}
\begin{Highlighting}[]
\NormalTok{|\textbackslash{}psi\_k|=\textbackslash{}frac\{A\_0\}\{N\}}
\end{Highlighting}
\end{Shaded}

同相整列時:

\begin{Shaded}
\begin{Highlighting}[]
\NormalTok{\textbackslash{}max|\textbackslash{}Psi\_N|=A\_0}
\end{Highlighting}
\end{Shaded}

\begin{center}\rule{0.5\linewidth}{0.5pt}\end{center}

\subsection{公理5:
個別位相不可読}\label{ux516cux74065-ux500bux5225ux4f4dux76f8ux4e0dux53efux8aad}

分類: 公理

\begin{Shaded}
\begin{Highlighting}[]
\NormalTok{\textbackslash{}theta\_k\textbackslash{}rightarrow\textbackslash{}theta\_k+\textbackslash{}alpha}
\end{Highlighting}
\end{Shaded}

\begin{center}\rule{0.5\linewidth}{0.5pt}\end{center}

\subsection{公理6:
位相差読出し}\label{ux516cux74066-ux4f4dux76f8ux5deeux8aadux51faux3057}

分類: 公理

\begin{Shaded}
\begin{Highlighting}[]
\NormalTok{\textbackslash{}Delta\_\{ij\}=\textbackslash{}theta\_i{-}\textbackslash{}theta\_j}
\end{Highlighting}
\end{Shaded}

読出し可能なペアチャンネル数:

\begin{Shaded}
\begin{Highlighting}[]
\NormalTok{\textbackslash{}binom\{N\}\{2\}=\textbackslash{}frac\{N(N{-}1)\}\{2\}}
\end{Highlighting}
\end{Shaded}

\begin{center}\rule{0.5\linewidth}{0.5pt}\end{center}

\subsection{公理7:
系同一視}\label{ux516cux74067-ux7cfbux540cux4e00ux8996}

分類: 公理

\begin{Shaded}
\begin{Highlighting}[]
\NormalTok{(N,A\_0,\textbackslash{}\{\textbackslash{}Delta\_\{ij\}\textbackslash{}\})}
\end{Highlighting}
\end{Shaded}

\begin{center}\rule{0.5\linewidth}{0.5pt}\end{center}

\subsection{公理8:
合成位相射影}\label{ux516cux74068-ux5408ux6210ux4f4dux76f8ux5c04ux5f71}

分類: 公理

\begin{Shaded}
\begin{Highlighting}[]
\NormalTok{S\_N=\textbackslash{}sum\_k e\^{}\{i\textbackslash{}theta\_k\}}
\end{Highlighting}
\end{Shaded}

\begin{Shaded}
\begin{Highlighting}[]
\NormalTok{\textbackslash{}Psi\_N=\textbackslash{}frac\{A\_0\}\{N\}S\_N}
\end{Highlighting}
\end{Shaded}

\begin{Shaded}
\begin{Highlighting}[]
\NormalTok{\textbackslash{}Theta\_N=\textbackslash{}arg\textbackslash{}Psi\_N}
\end{Highlighting}
\end{Shaded}

\begin{center}\rule{0.5\linewidth}{0.5pt}\end{center}

\subsection{公理9:
反転対称性非要請}\label{ux516cux74069-ux53cdux8ee2ux5bfeux79f0ux6027ux975eux8981ux8acb}

分類: 公理

反転対称性を要求しない。

\begin{center}\rule{0.5\linewidth}{0.5pt}\end{center}

\section{4. 外部読出し}\label{ux5916ux90e8ux8aadux51faux3057}

\subsection{公理10:
外部軸応答}\label{ux516cux740610-ux5916ux90e8ux8ef8ux5fdcux7b54}

分類: 公理候補

\begin{Shaded}
\begin{Highlighting}[]
\NormalTok{M\_\textbackslash{}alpha=e\^{}\{i\textbackslash{}alpha\}}
\end{Highlighting}
\end{Shaded}

\begin{Shaded}
\begin{Highlighting}[]
\NormalTok{Y(\textbackslash{}alpha)=\textbackslash{}frac\{A\_0\}\{N\}\textbackslash{}sum\_\{k=1\}\^{}\{N\}e\^{}\{i(\textbackslash{}theta\_k{-}\textbackslash{}alpha)\}}
\NormalTok{=e\^{}\{{-}i\textbackslash{}alpha\}\textbackslash{}Psi\_N}
\end{Highlighting}
\end{Shaded}

連続応答:

\begin{Shaded}
\begin{Highlighting}[]
\NormalTok{S\_N(\textbackslash{}alpha)=\textbackslash{}operatorname\{Re\}\textbackslash{}left[\textbackslash{}frac\{1\}\{N\}\textbackslash{}sum\_\{k=1\}\^{}\{N\}e\^{}\{i(\textbackslash{}theta\_k{-}\textbackslash{}alpha)\}\textbackslash{}right]}
\end{Highlighting}
\end{Shaded}

二値応答:

\begin{Shaded}
\begin{Highlighting}[]
\NormalTok{S\_N(\textbackslash{}alpha)=\textbackslash{}operatorname\{sgn\}\textbackslash{}operatorname\{Re\}\textbackslash{}left[\textbackslash{}frac\{1\}\{N\}\textbackslash{}sum\_\{k=1\}\^{}\{N\}e\^{}\{i(\textbackslash{}theta\_k{-}\textbackslash{}alpha)\}\textbackslash{}right]}
\end{Highlighting}
\end{Shaded}

\begin{center}\rule{0.5\linewidth}{0.5pt}\end{center}

\subsection{公理11:
外部読出し位相}\label{ux516cux740611-ux5916ux90e8ux8aadux51faux3057ux4f4dux76f8}

分類: 公理候補

\begin{Shaded}
\begin{Highlighting}[]
\NormalTok{\textbackslash{}Delta\textbackslash{}phi\_\textbackslash{}mu=\textbackslash{}phi\_\textbackslash{}mu\^{}\{(\textbackslash{}Psi)\}{-}\textbackslash{}phi\_\textbackslash{}mu\^{}\{(M)\}}
\end{Highlighting}
\end{Shaded}

\begin{Shaded}
\begin{Highlighting}[]
\NormalTok{\textbackslash{}mu=x,y,z,t}
\end{Highlighting}
\end{Shaded}

\begin{center}\rule{0.5\linewidth}{0.5pt}\end{center}

\subsection{公理12:
位置位相読出し}\label{ux516cux740612-ux4f4dux7f6eux4f4dux76f8ux8aadux51faux3057}

分類: 公理候補

\begin{Shaded}
\begin{Highlighting}[]
\NormalTok{x,y,z}
\end{Highlighting}
\end{Shaded}

\begin{center}\rule{0.5\linewidth}{0.5pt}\end{center}

\subsection{公理13:
時間位相読出し}\label{ux516cux740613-ux6642ux9593ux4f4dux76f8ux8aadux51faux3057}

分類: 公理候補

\begin{Shaded}
\begin{Highlighting}[]
\NormalTok{\textbackslash{}phi\_t\textbackslash{}equiv\textbackslash{}omega t\_\{\textbackslash{}mathrm\{read\}\}\textbackslash{}pmod\{2\textbackslash{}pi\}}
\end{Highlighting}
\end{Shaded}

\begin{center}\rule{0.5\linewidth}{0.5pt}\end{center}

\subsection{公理14:
被覆位相}\label{ux516cux740614-ux88abux8986ux4f4dux76f8}

分類: 公理候補

\begin{Shaded}
\begin{Highlighting}[]
\NormalTok{2π 周期読出し}
\NormalTok{4π 周期読出し}
\NormalTok{巻数差読出し}
\end{Highlighting}
\end{Shaded}

\begin{center}\rule{0.5\linewidth}{0.5pt}\end{center}

\section{5.
観測位相離散性}\label{ux89b3ux6e2cux4f4dux76f8ux96e2ux6563ux6027}

\subsection{公理15:
観測位相離散性}\label{ux516cux740615-ux89b3ux6e2cux4f4dux76f8ux96e2ux6563ux6027}

分類: 公理候補

\begin{Shaded}
\begin{Highlighting}[]
\NormalTok{\textbackslash{}sum\_\{\textbackslash{}mathrm\{cycle\}\}\textbackslash{}Delta\textbackslash{}phi=2\textbackslash{}pi m,\textbackslash{}qquad m\textbackslash{}in\textbackslash{}mathbb\{Z\}}
\end{Highlighting}
\end{Shaded}

\begin{Shaded}
\begin{Highlighting}[]
\NormalTok{W=\textbackslash{}frac\{1\}\{2\textbackslash{}pi\}\textbackslash{}oint d\textbackslash{}phi,\textbackslash{}qquad W\textbackslash{}in\textbackslash{}mathbb\{Z\}}
\end{Highlighting}
\end{Shaded}

読出しセル形式:

\begin{Shaded}
\begin{Highlighting}[]
\NormalTok{\textbackslash{}phi\_\{\textbackslash{}mathrm\{obs\}\}=m\textbackslash{}epsilon\_\textbackslash{}phi,\textbackslash{}qquad m\textbackslash{}in\textbackslash{}mathbb\{Z\}}
\end{Highlighting}
\end{Shaded}

位相面積候補:

\begin{Shaded}
\begin{Highlighting}[]
\NormalTok{Q\_\textbackslash{}phi=\textbackslash{}sum\_\{i\textless{}j\}\textbackslash{}sin\^{}2\textbackslash{}frac\{\textbackslash{}Delta\_\{ij\}\}\{2\}}
\end{Highlighting}
\end{Shaded}

\begin{center}\rule{0.5\linewidth}{0.5pt}\end{center}

\section{6. 波形・局在}\label{ux6ce2ux5f62ux5c40ux5728}

\subsection{作業公理A:
状態波形}\label{ux4f5cux696dux516cux7406a-ux72b6ux614bux6ce2ux5f62}

分類: 作業仮説

\begin{Shaded}
\begin{Highlighting}[]
\NormalTok{A(\textbackslash{}theta)=\textbackslash{}sum\_n A\_n\textbackslash{}cos(n\textbackslash{}theta+\textbackslash{}phi\_n)}
\end{Highlighting}
\end{Shaded}

\begin{Shaded}
\begin{Highlighting}[]
\NormalTok{Z(\textbackslash{}theta)=\textbackslash{}sum\_n r\_n e\^{}\{i(n\textbackslash{}theta+\textbackslash{}phi\_n)\}}
\end{Highlighting}
\end{Shaded}

\begin{center}\rule{0.5\linewidth}{0.5pt}\end{center}

\subsection{定義A1:
正規化等振幅奇数倍音局在波}\label{ux5b9aux7fa9a1-ux6b63ux898fux5316ux7b49ux632fux5e45ux5947ux6570ux500dux97f3ux5c40ux5728ux6ce2}

分類: 定義

\begin{Shaded}
\begin{Highlighting}[]
\NormalTok{N\_h=2K{-}1}
\end{Highlighting}
\end{Shaded}

\begin{Shaded}
\begin{Highlighting}[]
\NormalTok{S\_\{N\_h\}(u)=\textbackslash{}frac\{1\}\{K\}\textbackslash{}sum\_\{m=0\}\^{}\{K{-}1\}\textbackslash{}cos((2m+1)u)}
\end{Highlighting}
\end{Shaded}

\begin{Shaded}
\begin{Highlighting}[]
\NormalTok{Z\_\{N\_h\}(u)=\textbackslash{}frac\{1\}\{K\}\textbackslash{}sum\_\{m=0\}\^{}\{K{-}1\}e\^{}\{i(2m+1)u\}}
\end{Highlighting}
\end{Shaded}

\begin{Shaded}
\begin{Highlighting}[]
\NormalTok{\textbackslash{}Psi\_\{N\_h\}(u)=A S\_\{N\_h\}(u)}
\end{Highlighting}
\end{Shaded}

\begin{Shaded}
\begin{Highlighting}[]
\NormalTok{\textbackslash{}Psi\_\{N\_h\}(u)=A Z\_\{N\_h\}(u)}
\end{Highlighting}
\end{Shaded}

\begin{Shaded}
\begin{Highlighting}[]
\NormalTok{\textbackslash{}frac\{A\}\{K\}}
\end{Highlighting}
\end{Shaded}

\begin{center}\rule{0.5\linewidth}{0.5pt}\end{center}

\subsection{定理A1:
代表振幅保存}\label{ux5b9aux7406a1-ux4ee3ux8868ux632fux5e45ux4fddux5b58}

分類: 導出済み定理

\begin{Shaded}
\begin{Highlighting}[]
\NormalTok{S\_\{N\_h\}(0)=1}
\end{Highlighting}
\end{Shaded}

\begin{Shaded}
\begin{Highlighting}[]
\NormalTok{\textbackslash{}Psi\_\{N\_h\}(0)=A}
\end{Highlighting}
\end{Shaded}

\begin{center}\rule{0.5\linewidth}{0.5pt}\end{center}

\subsection{定理A2:
最高奇数倍音次数と局在幅}\label{ux5b9aux7406a2-ux6700ux9ad8ux5947ux6570ux500dux97f3ux6b21ux6570ux3068ux5c40ux5728ux5e45}

分類: 導出済み定理

\begin{Shaded}
\begin{Highlighting}[]
\NormalTok{K=\textbackslash{}frac\{N\_h+1\}\{2\}}
\end{Highlighting}
\end{Shaded}

\begin{Shaded}
\begin{Highlighting}[]
\NormalTok{\textbackslash{}lambda\_\{N\_h\}=\textbackslash{}frac\{\textbackslash{}lambda\_0\}\{N\_h\}}
\end{Highlighting}
\end{Shaded}

\begin{Shaded}
\begin{Highlighting}[]
\NormalTok{\textbackslash{}Delta x\_\{N\_h\}\textbackslash{}propto\textbackslash{}frac\{\textbackslash{}lambda\_0\}\{N\_h\}}
\end{Highlighting}
\end{Shaded}

\begin{Shaded}
\begin{Highlighting}[]
\NormalTok{N\_\{h,\textbackslash{}rm req\}\textbackslash{}sim\textbackslash{}frac\{\textbackslash{}lambda\_0\}\{\textbackslash{}Delta x\}}
\end{Highlighting}
\end{Shaded}

\begin{Shaded}
\begin{Highlighting}[]
\NormalTok{N\_\{h,\textbackslash{}rm req\}\^{}\{(t)\}\textbackslash{}sim\textbackslash{}frac\{T\_0\}\{\textbackslash{}Delta t\}}
\end{Highlighting}
\end{Shaded}

\begin{center}\rule{0.5\linewidth}{0.5pt}\end{center}

\subsection{帰結A1:
空間・時間局在の同一構成}\label{ux5e30ux7d50a1-ux7a7aux9593ux6642ux9593ux5c40ux5728ux306eux540cux4e00ux69cbux6210}

分類: 導出済み帰結

\begin{Shaded}
\begin{Highlighting}[]
\NormalTok{\textbackslash{}Psi(\textbackslash{}chi,\textbackslash{}tau)=A S\_\{N\_\{h,\textbackslash{}chi\}\}(\textbackslash{}chi{-}\textbackslash{}chi\_0)S\_\{N\_\{h,\textbackslash{}tau\}\}(\textbackslash{}tau{-}\textbackslash{}tau\_0)e\^{}\{i\textbackslash{}phi\}}
\end{Highlighting}
\end{Shaded}

\begin{center}\rule{0.5\linewidth}{0.5pt}\end{center}

\subsection{作業公理B:
局在}\label{ux4f5cux696dux516cux7406b-ux5c40ux5728}

分類: 作業仮説

\begin{Shaded}
\begin{Highlighting}[]
\NormalTok{S\_\{N\_h\}(u)=\textbackslash{}frac\{1\}\{K\}\textbackslash{}sum\_\{m=0\}\^{}\{K{-}1\}\textbackslash{}cos((2m+1)u)}
\end{Highlighting}
\end{Shaded}

\begin{Shaded}
\begin{Highlighting}[]
\NormalTok{Z\_\{N\_h\}(u)=\textbackslash{}frac\{1\}\{K\}\textbackslash{}sum\_\{m=0\}\^{}\{K{-}1\}e\^{}\{i(2m+1)u\}}
\end{Highlighting}
\end{Shaded}

\begin{center}\rule{0.5\linewidth}{0.5pt}\end{center}

\subsection{作業公理C:
結合入口}\label{ux4f5cux696dux516cux7406c-ux7d50ux5408ux5165ux53e3}

分類: 作業仮説

相互作用の入口は低次ベース位相の整合である。

\begin{center}\rule{0.5\linewidth}{0.5pt}\end{center}

\subsection{作業公理D:
観測写像}\label{ux4f5cux696dux516cux7406d-ux89b3ux6e2cux5199ux50cf}

分類: 作業仮説

\begin{Shaded}
\begin{Highlighting}[]
\NormalTok{s\_\textbackslash{}alpha=\textbackslash{}mathrm\{Loc\}\textbackslash{}left[\textbackslash{}int A(\textbackslash{}theta)M\_\textbackslash{}alpha(\textbackslash{}theta)d\textbackslash{}theta\textbackslash{}right]}
\end{Highlighting}
\end{Shaded}

\begin{center}\rule{0.5\linewidth}{0.5pt}\end{center}

\subsection{作業公理E:
相関強度二乗}\label{ux4f5cux696dux516cux7406e-ux76f8ux95a2ux5f37ux5ea6ux4e8cux4e57}

分類: 作業仮説

\begin{Shaded}
\begin{Highlighting}[]
\NormalTok{P(\textbackslash{}alpha)=\textbackslash{}frac\{|C\_\textbackslash{}alpha|\^{}2\}\{\textbackslash{}sum\_\textbackslash{}beta |C\_\textbackslash{}beta|\^{}2\}}
\end{Highlighting}
\end{Shaded}

\begin{center}\rule{0.5\linewidth}{0.5pt}\end{center}

\subsection{作業公理F:
低次ベース同期共有}\label{ux4f5cux696dux516cux7406f-ux4f4eux6b21ux30d9ux30fcux30b9ux540cux671fux5171ux6709}

分類: 作業仮説

\begin{Shaded}
\begin{Highlighting}[]
\NormalTok{\textbackslash{}theta\_A{-}\textbackslash{}theta\_B=\textbackslash{}Delta}
\end{Highlighting}
\end{Shaded}

\begin{center}\rule{0.5\linewidth}{0.5pt}\end{center}

\subsection{作業公理G:
同期復調不能化}\label{ux4f5cux696dux516cux7406g-ux540cux671fux5fa9ux8abfux4e0dux80fdux5316}

分類: 作業仮説

\begin{Shaded}
\begin{Highlighting}[]
\NormalTok{V(t)=\textbackslash{}left|\textbackslash{}sum\_n w\_n e\^{}\{i\textbackslash{}epsilon\_n(t)\}\textbackslash{}right|}
\end{Highlighting}
\end{Shaded}

\begin{Shaded}
\begin{Highlighting}[]
\NormalTok{V(t)=\textbackslash{}left|\textbackslash{}sum\_n w\_n e\^{}\{{-}\textbackslash{}frac12\textbackslash{}sigma\_n\^{}2(t)\}\textbackslash{}right|}
\end{Highlighting}
\end{Shaded}

\begin{center}\rule{0.5\linewidth}{0.5pt}\end{center}

\subsection{作業公理H:
階層性}\label{ux4f5cux696dux516cux7406h-ux968eux5c64ux6027}

分類: 作業仮説

\begin{Shaded}
\begin{Highlighting}[]
\NormalTok{低次ベースモード}
\NormalTok{+ 高次局在モード}
\end{Highlighting}
\end{Shaded}

\begin{center}\rule{0.5\linewidth}{0.5pt}\end{center}

\section{7.
読出し・記憶・複素表示}\label{ux8aadux51faux3057ux8a18ux61b6ux8907ux7d20ux8868ux793a}

\subsection{I/Q 読出し}\label{iq-ux8aadux51faux3057}

分類: 理論的知見

\begin{Shaded}
\begin{Highlighting}[]
\NormalTok{z=a+ib}
\end{Highlighting}
\end{Shaded}

\begin{Shaded}
\begin{Highlighting}[]
\NormalTok{a=r\textbackslash{}cos\textbackslash{}theta,\textbackslash{}qquad b=r\textbackslash{}sin\textbackslash{}theta}
\end{Highlighting}
\end{Shaded}

\begin{Shaded}
\begin{Highlighting}[]
\NormalTok{u=r\textbackslash{}cos\textbackslash{}theta+r\textbackslash{}sin\textbackslash{}theta}
\NormalTok{=\textbackslash{}sqrt2 r\textbackslash{}cos\textbackslash{}left(\textbackslash{}theta{-}\textbackslash{}frac\{\textbackslash{}pi\}\{4\}\textbackslash{}right)}
\end{Highlighting}
\end{Shaded}

\begin{center}\rule{0.5\linewidth}{0.5pt}\end{center}

\subsection{倍音情報読出し}\label{ux500dux97f3ux60c5ux5831ux8aadux51faux3057}

分類: 理論的知見

\begin{Shaded}
\begin{Highlighting}[]
\NormalTok{Z(\textbackslash{}theta)=\textbackslash{}sum\_n r\_n e\^{}\{i(n\textbackslash{}theta+\textbackslash{}phi\_n)\}}
\end{Highlighting}
\end{Shaded}

\begin{Shaded}
\begin{Highlighting}[]
\NormalTok{P(Z)=\textbackslash{}operatorname\{Re\}(Z)+\textbackslash{}operatorname\{Im\}(Z)}
\end{Highlighting}
\end{Shaded}

\begin{Shaded}
\begin{Highlighting}[]
\NormalTok{P(Z)=\textbackslash{}sum\_n \textbackslash{}sqrt2 r\_n}
\NormalTok{\textbackslash{}cos\textbackslash{}left(n\textbackslash{}theta+\textbackslash{}phi\_n{-}\textbackslash{}frac\{\textbackslash{}pi\}\{4\}\textbackslash{}right)}
\end{Highlighting}
\end{Shaded}

\begin{center}\rule{0.5\linewidth}{0.5pt}\end{center}

\subsection{保存量側と読出し側}\label{ux4fddux5b58ux91cfux5074ux3068ux8aadux51faux3057ux5074}

分類: 理論的知見

\begin{Shaded}
\begin{Highlighting}[]
\NormalTok{z\textbackslash{}bar z=a\^{}2+b\^{}2}
\end{Highlighting}
\end{Shaded}

\begin{Shaded}
\begin{Highlighting}[]
\NormalTok{z\^{}2=a\^{}2{-}b\^{}2+2iab}
\end{Highlighting}
\end{Shaded}

\begin{center}\rule{0.5\linewidth}{0.5pt}\end{center}

\subsection{二乗レジストリ}\label{ux4e8cux4e57ux30ecux30b8ux30b9ux30c8ux30ea}

分類: 作業仮説

\begin{Shaded}
\begin{Highlighting}[]
\NormalTok{\textbackslash{}mathcal\{A\}\_0=\textbackslash{}sum\_k A\_k\^{}2}
\end{Highlighting}
\end{Shaded}

\begin{Shaded}
\begin{Highlighting}[]
\NormalTok{a\_0=\textbackslash{}sum\_k |A\_k|\^{}2=\textbackslash{}sum\_k A\_k\textbackslash{}bar A\_k}
\end{Highlighting}
\end{Shaded}

\begin{center}\rule{0.5\linewidth}{0.5pt}\end{center}

\section{8. 統計分類}\label{ux7d71ux8a08ux5206ux985e}

\subsection{ペア位相差}\label{ux30daux30a2ux4f4dux76f8ux5dee}

分類: 定義

\begin{Shaded}
\begin{Highlighting}[]
\NormalTok{\textbackslash{}Delta\_\{ij\}=\textbackslash{}theta\_i{-}\textbackslash{}theta\_j}
\end{Highlighting}
\end{Shaded}

\begin{Shaded}
\begin{Highlighting}[]
\NormalTok{\textbackslash{}binom\{N\}\{2\}=\textbackslash{}frac\{N(N{-}1)\}\{2\}}
\end{Highlighting}
\end{Shaded}

ペア合成:

\begin{Shaded}
\begin{Highlighting}[]
\NormalTok{\textbackslash{}psi\_i+\textbackslash{}psi\_j}
\NormalTok{=2A\textbackslash{}cos\textbackslash{}frac\{\textbackslash{}Delta\_\{ij\}\}\{2\}e\^{}\{i(\textbackslash{}theta\_i+\textbackslash{}theta\_j)/2\}}
\end{Highlighting}
\end{Shaded}

ペア干渉強度:

\begin{Shaded}
\begin{Highlighting}[]
\NormalTok{I\_\{ij\}=4A\^{}2\textbackslash{}cos\^{}2\textbackslash{}frac\{\textbackslash{}Delta\_\{ij\}\}\{2\}}
\end{Highlighting}
\end{Shaded}

\begin{center}\rule{0.5\linewidth}{0.5pt}\end{center}

\subsection{重なり度}\label{ux91cdux306aux308aux5ea6}

分類: 定義

\begin{Shaded}
\begin{Highlighting}[]
\NormalTok{R\_N=\textbackslash{}left|\textbackslash{}frac\{A\_0\}\{N\}\textbackslash{}sum\_\{k=1\}\^{}\{N\}e\^{}\{i\textbackslash{}theta\_k\}\textbackslash{}right|}
\end{Highlighting}
\end{Shaded}

\begin{Shaded}
\begin{Highlighting}[]
\NormalTok{\textbackslash{}rho\_N=\textbackslash{}frac\{R\_N\}\{A\_0\}}
\NormalTok{=\textbackslash{}frac\{1\}\{N\}\textbackslash{}left|\textbackslash{}sum\_\{k=1\}\^{}\{N\}e\^{}\{i\textbackslash{}theta\_k\}\textbackslash{}right|}
\end{Highlighting}
\end{Shaded}

\begin{Shaded}
\begin{Highlighting}[]
\NormalTok{0\textbackslash{}le\textbackslash{}rho\_N\textbackslash{}le1}
\end{Highlighting}
\end{Shaded}

\begin{center}\rule{0.5\linewidth}{0.5pt}\end{center}

\subsection{フェルミオン度}\label{ux30d5ux30a7ux30ebux30dfux30aaux30f3ux5ea6}

分類: 定義

\begin{Shaded}
\begin{Highlighting}[]
\NormalTok{F\_N=}
\NormalTok{\textbackslash{}frac\{2\}\{N(N{-}1)\}}
\NormalTok{\textbackslash{}sum\_\{i\textless{}j\}\textbackslash{}sin\^{}2\textbackslash{}frac\{\textbackslash{}Delta\_\{ij\}\}\{2\}}
\end{Highlighting}
\end{Shaded}

\begin{Shaded}
\begin{Highlighting}[]
\NormalTok{B\_N=1{-}F\_N}
\end{Highlighting}
\end{Shaded}

\begin{Shaded}
\begin{Highlighting}[]
\NormalTok{B\_N=}
\NormalTok{\textbackslash{}frac\{2\}\{N(N{-}1)\}}
\NormalTok{\textbackslash{}sum\_\{i\textless{}j\}\textbackslash{}cos\^{}2\textbackslash{}frac\{\textbackslash{}Delta\_\{ij\}\}\{2\}}
\end{Highlighting}
\end{Shaded}

\begin{Shaded}
\begin{Highlighting}[]
\NormalTok{F\_N=\textbackslash{}frac\{N\}\{2(N{-}1)\}(1{-}\textbackslash{}rho\_N\^{}2)}
\end{Highlighting}
\end{Shaded}


\end{document}
